\subsection{Transfer-only}

The transfer-only setting is generally used so as to not burn armies in case your opponent captures a territory. For example in this \href{https://www.warzone.com/MultiPlayer?GameID=40619054}{game} on Turn 8 I send in my first move with transfer only those 2 troops as I want to make sure they don't go attacking. Or similarly in \href{https://www.warzone.com/MultiPlayer?GameID=39911346}{this}, Aika on turn 16 sends his armies forward north with transfer only in case I haven't attacked and captured. Same for \href{https://www.warzone.com/MultiPlayer?GameID=41380732}{here}, where ps on Turn 19 uses transfer-only on his first move.
\\
\\
The more unique/niche idea behind transfer-only is that if it fails, it locks your armies in place, which creates new opportunities. For example in this \href{https://www.warzone.com/MultiPlayer?GameID=40836522}{Turkey LD} game, Rufus expects me to first move him, so he performs a 2army transfer-only move to a neutral that will always fail, effectively locking 2 movable armies in place. This is for the scenario I first move him with 6, leaving him with 2 movable armies. If I don't attack him on my first move, his armies from north have joined in and his last move completion of the bonus will succeed. If I attack him, nothing happens and he retains his army advantage. Similarly, on this \href{https://www.warzone.com/MultiPlayer?GameID=34661130}{game}, on Turn 6, adrenaline wants to retain his stack in place at Zama, in Africa. So he does a 12-army transfer-only move to a neutral followed by a 3-army attack
\\
\\
The fact that failed transfer-only moves result in locked armies can be further used to optimize the usage of the Blockade card, as shown in the Blockade section.